\scanpage{29}%
\mspage{14}\space
$2$ nilpotent, which makes it possible to put the curve in the form (4.1)
\[
y^2+xy=x^3+a_2\,x+a_6
\]
and an isomorphism of such a curve onto another is described by $s$ satisfying (4.3) for $u=1+2s$. When $j=j'$, the first of these equations
\[
(1+2s)^2\,a'_2=a_2-s-s^2
\]
defines an étale covering of degree two, and it implies the second ; to see this, one reduces to assuming $a_2=a'_2$ and observes that by (4.4)
\[
1/j=a_6+2\,R(a_2,a_6)
\]
for a rational function $R$ .

To verify (II) in the cases not covered by (III), one may suppose that $S$ is the spectrum of an algebraically closed field and that $c_4=0$ or $c_6=0$ \tikz[baseline=(punct.base)]{\node[inner sep=0pt] (punct) {,};\node[inner sep=0pt,anchor=north] at ([yshift=-0.16em]punct.south) {.};} i.e. that $j=0$ or $j=1728$. Let $p$ be the characteristic of the field,

a) $p\ne 2,3$ : one puts the curve in the form (2.1) ; if $c_4=0$ (resp. $c_6=0$) one will have $a_4=0$ (resp. $a_6=0$) and one can by homothety transform the curve into the standard curve
\[
y^2=x^3+1 \qquad (\text{if }c_4=0) \tag*{(6.5)}
\]
\[
y^2=x^3-x \qquad (\text{if }c_6=0) \tag*{(6.6)}
\]

b) $\unsure{\mathbb{\text{\makebox[0pt][l]{\kern.45pt$p$}$p$}}|p}=3$ : one puts the curve in the form (6.1) ; solving (6.2) for $a'_4=1$ and $a'_6=0$, one puts the curve in the standard form%
% archive sheet number: 28
